% !TeX root = ../main.tex
% apx_b_static_scope.tex -- one section of Appendix B, kept in its own file so that a SINGLE
% commented \input line in apx_b_errata.tex suppresses it (author's condition, PENDENCIAS_RESOLVIDOS 2.4 (arquivado 2026-07-30), was 2.2 before the 2026-07-29 renumber:
% "deixe isso facil de ser comentado"). Provenance for every number is in the comment block
% at that \input site. Do not inline this file back into the appendix.

\section{The scope of the static task in Chapter~\ref{ch:courb}}
\label{apx:errata:static-scope}

Chapter~\ref{ch:courb} reports two tasks. One is sequential: predict the category of the
next place a user visits, from the history of that user's check-ins. The other is static:
classify the category of a place from that place's own representation. This section
records what we established after publication about the second one, and it bears on how
that chapter's category results should be read.

% [2026-07-30, author decision in PENDENCIAS_RESOLVIDOS 2.4 (arquivado 2026-07-30)] REWRITTEN FOR CLARITY at his request: "a
% explicação não está bem didatica e clara, requere bastante esforço para entender e varias
% leituras". His example was the exact sentence this replaces -- "built from a fine-grained class
% label attached to each place" -- which never tells the reader that the thing in question is the
% VENUE'S OWN NAME. It now says so in the second sentence, with three real values.
%
% THE RANGE 284-365 IS CORRECT AND WAS RE-CONFIRMED. Do not "fix" it to 377: I did exactly that on
% 2026-07-30 and it was wrong, because the number depends on WHICH COLUMN and AT WHAT GRANULARITY.
% The pipeline's fine class is the parquet's `spot` column, which
% research/embeddings/hgi/preprocess.py:62 renames to `fclass`, then drops rows with a null category
% (:64) and reduces to ONE ROW PER placeid (:75-80). Reproduced on the five state parquets:
%   Alabama 284 | Arizona 305 | Florida 324 | California 333 | Texas 365   -> the range in the text
%   values mapping to more than one top-level category: 0, in every state
% The parquet ALSO has a `spot_categories` column, which is a per-check-in JSON string. Counting
% distinct values there, over raw check-in rows without the placeid dedup, gives 284/306/325/334/377
% -- a different quantity that has no role in the pipeline. `fclass` is not a stale name: it is what
% poi2vec.py requires and what the embedding is actually keyed on.
% The example strings are verbatim values of that column: 'Coffee Shop', 'Seafood', 'Airport'.
The problem is easiest to see by asking what the static task is actually given. Each place in the
Gowalla data carries a fine-grained venue type, a string such as \texttt{Coffee Shop},
\texttt{Seafood}, or \texttt{Airport}. The place embedding that Chapter~\ref{ch:courb} feeds to the
static task is a lookup table on that string: two places of the same venue type receive the same
vector. The target the task must predict is the top-level category, and \texttt{Coffee Shop} belongs
to Food no matter which coffee shop it is.

The input therefore contains the answer. We confirmed this holds without exception across the five
Gowalla state subsets the collection uses: the venue type takes between 284 and 365 distinct values
per state, and in every state each of those values maps to exactly one of the seven top-level
categories. Not one is ambiguous. A model given the venue type and asked for the category is being
asked to reproduce a lookup it already holds, so the reported accuracy on the static task measures
that mapping rather than any learned semantic inference.

Three qualifications matter, and none of them is a softening.

The first is that the sequential task is unaffected. Its input is a history of visits and
its target is a category the model has not seen for that step, so no equivalent identity
exists between input and label. Every claim Chapter~\ref{ch:courb} makes about the
sequential task, including the improvement that the decomposed representation produces,
stands as published.

The second is that Chapter~\ref{ch:cbic} has a version of the same problem, milder but present, and
saying so is more honest than drawing a line between the two. That chapter builds each place's input
feature from the average of its neighbors' categories and deliberately leaves the place itself out.
We verified that this exclusion holds exactly: a place's own category vector is absent from its own
input feature. The exclusion does not close the channel, because the input feature is not what the
static task reads. Each neighbor's feature vector is itself an average over that neighbor's own
neighborhood, and the place belongs to those neighborhoods, so one round of graph attention returns
the place's category to its embedding at a measured mean weight of 0.10 against a total own-category
weight of 0.39. Removing that contribution and leaving everything else unchanged lowers a probe of
the place's own category from 0.46 to 0.30 macro-F1, against a majority-class floor of 0.07.
Relabeling a single place, which never alters that place's own input feature, changes its embedding
for every place that receives at least one message. The difference between the two chapters is
therefore one of degree: an exact lookup in Chapter~\ref{ch:courb}, a diluted average recovered
through one hop in Chapter~\ref{ch:cbic}.
% [NEEDS SIGN-OFF 46 | OPEN] Author asked for an independent code-level verification rather than
% accepting either his own belief (no leak, since the input feature excludes the place's own label)
% or the prior round's finding (D-01) at face value. A fresh audit (see src_utils/, DGI-leak-audit)
% traced the executed path with the real encoder on real Florida/California data and confirmed a
% leak, quantified above. It also corrected D-01 itself: the encoder is GATConv
% (research/embeddings/dgi/model/DGIModule.py:12), not GCNConv(cached=True) (an unused sibling
% module); and the reentry route is not edge_index symmetry (edge_index is directed and asymmetric,
% verified on the real preprocessed graphs) but the undirected neighbor-mean construction of each
% neighbor's own input feature. A causal intervention (relabeling one place and re-running the real
% encoder, with that place's own input row unchanged) confirmed the channel is not an artifact of the
% linear-algebra decomposition. Full trace, code, and figures: SOURCE_LEDGER_dgi_selfleak.md and
% dgi_selfleak_audit.png (this session's artifacts).
% THIS PARAGRAPH NOW SAYS SOMETHING STRONGER ABOUT CHAPTER 3 THAN THE AUTHOR'S ORIGINAL RULING
% ASSUMED (he believed the input feature's exclusion closed the channel; the measurement says it
% does not), and the audit's own author explicitly flagged that the flag should stay until he reads
% it. Two open questions for him: (1) whether the +0.165 macro-F1 delta changes how Chapter 3's
% static-task numbers should be reported at all, beyond this appendix; (2) an incidental finding in
% the audit's ledger \S F, outside the leak question: the DGI contrastive objective as implemented
% appears degenerate (positive and negative score sets are identical as multisets; measured loss
% floor matches 2*ln2 to six decimals), which may mean "trained DGI embedding" does not describe
% what Chapter 3 actually used. Not acted on here; his call.

The third is that Chapter~\ref{ch:mobiwac} does not inherit the problem, for the reason
that chapter gives on its own terms: it replaces the place-level embedding with a
check-in-level representation, whose input is a single visit, and the identity described
above does not arise.

The statement belongs here rather than in the chapter because Chapter~\ref{ch:courb} is a published,
co-authored article, and this appendix is where this collection records the difference between what
a published text says and what we know now. The
measurement is recorded in the repository released with this work.
